\section{数列怎样收敛}
\label{sec:02-Calculus/05-Sequences-and-Series/01}

\subsection{先有困惑，后有定义}

碰到这样一个递推数列：
\[
a_1=0,\qquad a_{n+1}=\sqrt{2+a_n}.
\]

先算几项感受一下：\(a_1=0\)，\(a_2=\sqrt2\approx1.414\)，\(a_3=\sqrt{2+\sqrt2}\approx1.848\)，\(a_4\approx1.962\)，\(a_5\approx1.990\)……数字似乎在往 2 靠拢。感觉不错。

但问题来了：你怎么知道它一定会收敛，而不是在 2 附近无限次地贴上去又弹开？你怎么知道它不是慢慢地攀升到比如 2.000137 然后停在那里？算五六项看到趋势就能下结论吗？

这是所有收敛概念要回答的根本问题。``越来越近''不是一个可靠的数学判断——序列 \(a_n=2+(1/n)\sin(n)\) 在 \(n\) 很大时每隔一段时间就会漂离 2 一小截，然后又回来。所以需要比``看起来近了''更精确的说法。

\subsection{把``最终全部进入窗口''固定为定义}

数列 \((a_n)\) 就是把每个自然数 \(n\) 映射到一个实数 \(a_n\) 的函数：\(a:\N\to\R\)。顺序固定，重复也保留——\(\{1,-1,1,-1,\dots\}\) 不是集合 \(\{1,-1\}\)。

我们说 \(a_n\to L\)，精确含义是：
\[
\forall\varepsilon>0,\;\exists N\in\N,\;\forall n\ge N,\;|a_n-L|<\varepsilon.
\]

翻译成日常语言：无论你提出多窄的容忍窗口（\(\varepsilon\)），从某一项（\(N\)）开始，之后所有的项都老老实实待在窗口里面。三个关键点：
\begin{itemize}
\tightlist
\item ``从某一项之后的所有项''——允许前面有限项任意乱跑；
\item ``所有''——不允许尾部有无穷多项反复逃出窗口；
\item ``无论多窄''——\(\varepsilon\) 取任意正数都能找到对应的 \(N\)。
\end{itemize}

验证经典例子 \(1/n\to0\)：给定 \(\varepsilon>0\)，需要找到 \(N\) 使得 \(n\ge N\) 时 \(1/n<\varepsilon\)。取 \(N\) 为大于 \(1/\varepsilon\) 的任意自然数即可——Archimedes 性质保证这样的自然数存在。

对照 \((-1)^n\)：它对任何候选极限 \(L\) 都能被证伪。假设 \(L=1\)，取 \(\varepsilon=0.5\)，则奇数项 \(a_{2k+1}=-1\) 与 1 的距离为 2，远大于 0.5，而且有无限多项。无论把 \(N\) 设得多大，总还有奇数项在窗口外面。同样可以反驳 \(L=-1\) 或 \(L\) 为任何其他值。\((-1)^n\) 不收敛。

\subsection{收敛序列的代数性质}

极限若存在，必唯一。假如一个序列有两个不同的极限 \(L_1\neq L_2\)，取 \(\varepsilon=|L_1-L_2|/3\)，尾部的项不可能同时离两个极限都足够近——三角形不等式会让这个美梦落空。

收敛序列必须有界：尾部都在极限的邻域内，有限项只有有限个，取最大绝对值就得到一个界。注意反过来不对：有界序列未必收敛，\((-1)^n\) 就是反例——它的值被约束在 \(-1\) 和 1 之间，但它不走近任何一个固定数。

四则运算上，若 \(a_n\to A\)、\(b_n\to B\)，则 \(a_n+b_n\to A+B\)，\(a_n b_n\to AB\)，且当 \(B\neq0\) 时 \(a_n/b_n\to A/B\)（分母从某项起不为零）。这些性质让你可以把复杂序列拆成已知极限的组合。

\subsection{单调有界：不需要知道极限也知道它存在}

有界不一定收敛，单调也不一定收敛（\(a_n=n\) 单调递增却无界）。但两个条件合在一起刚好够用：

单调递增且有上界的实数列一定收敛，极限就是所有项的最小上界（上确界）。递减情形对称。

这个定理的威力在于：它保证极限存在，但不需要你事先知道极限是多少。它依赖实数系的完备性——有理数序列 \(\sqrt2\) 的十进制逼近有上界且单调，但在有理数范围内没有极限；实数域填上了所有的``空隙''。

回到开头那个递推数列。用单调有界定理分两步验证收敛性：

第一步：证明有界。先猜上界为 2。\(a_1=0\le2\)；假设 \(a_n\le2\)，则
\[
a_{n+1}=\sqrt{2+a_n}\le\sqrt{2+2}=2,
\]
归纳封口。下界 \(a_n\ge0\) 同样归纳可得。

第二步：证明单调。在 \(0\le a_n\le2\) 的条件下比较 \(a_{n+1}\) 和 \(a_n\)：
\[
a_{n+1}\ge a_n
\Longleftrightarrow
\sqrt{2+a_n}\ge a_n
\Longleftrightarrow
2+a_n\ge a_n^2.
\]
后者等价于 \((a_n-2)(a_n+1)\le0\)，而 \(a_n\in[0,2]\) 恰好满足。所以 \(a_{n+1}\ge a_n\)，数列单调不减。

有上界 + 单调递增 \(\Longrightarrow\) 收敛。到这一步，极限的存在性已经板上钉钉。

\subsection{先证存在，才能代入极限}

很多人在处理递归数列时上来就写：``令 \(n\to\infty\)，设极限为 \(L\)，则 \(L=\sqrt{2+L}\)，解得 \(L=2\) 或 \(L=-1\)。''然后因为别的原因排除 \(-1\)，宣布极限为 2。

这个推导在逻辑上颠倒了。\(L=\sqrt{2+L}\) 这一步做了两件事：一是假设极限存在，二是假设递推式中的运算可以在极限号下保持。这两步都需要前提——你还没证明极限存在，怎么能把它代入等式？

正确的顺序是：先用单调有界定理证明极限存在；然后因为平方根函数连续，可以把极限取进递推式内部，得到 \(L=\sqrt{2+L}\)；由 \(a_n\ge0\) 排除 \(L=-1\)，得 \(L=2\)。

两种顺序写出来的答案一样，但第一种是跳过了关键步骤，第二种才建立了从定义到结论的完整链条。

\subsection{子列：从振荡中提取证据}

从原序列中按递增的下标序列 \(n_1<n_2<n_3<\cdots\) 取出 \(a_{n_1},a_{n_2},a_{n_3},\ldots\)，称为子列。若原序列收敛于 \(L\)，则每个子列也收敛于 \(L\)——因为尾部全部靠近 \(L\)，子列的尾部当然也不例 外。

这个方向最常用的方式是对它取逆否：若能找到两个子列趋向不同极限，原序列一定不收敛。对 \((-1)^n\)，偶数项子列 \(a_{2k}=1\to1\)，奇数项子列 \(a_{2k+1}=-1\to-1\)，两个极限不同，所以原序列发散。这个论证把``振荡''从直觉图像变成了逻辑上无懈可击的证明。

方向相反的结论是 Bolzano--Weierstrass 定理：每个有界实数序列至少包含一个收敛子列。这个定理不要求原序列收敛——它说的是，即使序列整体在区间内乱跳，你总能在其中抽出某种趋于稳定的趋势。证明的核心是区间套：把有界区间反复二分，至少有一半包含无穷多项，不断缩小区间套出一个极限点，再逐项抽出收敛到该点的子列。

这个定理在后面很多地方都会出现：紧致性论证、极值存在定理、Cauchy 列的处理，都离不开``总有一个收敛子列''这个保障。它告诉你，有界的混乱不是彻底的混乱——总有一丝秩序藏在里面。

\subsection{定义没告诉你的事}

收敛的定义告诉你极限在哪里，但完全不涉及收敛的速度。\(\lim 1/n=0\) 和 \(\lim 1/\ln(n)=0\) 都收敛于同一个数，但前者每 10 倍 \(n\) 精确一位小数，后者每 10 倍 \(n\) 只能前进 \(1/\ln10\) 左右。定义不区分这些。

单调有界定理保证极限存在，但它是一个纯粹的存在性结论。很多递归定义的数列，极限只知在某个范围内，却写不出初等闭式——数值计算和不等式估计才是真正的工作。
